\section{Bluetooth Low Energy: standard overview and throughput analysis}
The choice of radios and communication protocol is always a hard one when dealing with sensor networks. As different technologies offer different performances in different areas, the choice should be coherent with the kind of data we are dealing with. For example, some aspects that should be considered are throughput, latency, number of nodes and compatibility: while the first two could be critical in healthcare applications, compatibility could be crucial when dealing with sensor devices aimed at the mass market. In particular, the Bluetooth 4.0 standard (also called  Bluetooth Low Energy) aims at maximizing device connectivity and energy efficiency, with the drawback of a limited throughput.

\paragraph{The Bluetooth Application Layer.} When dealing with BLE chips, usually the final user can have control on the Application Layer of the communication protocol, while lower levels (Host and Controller) usually can't be optimized by the developer. In particular, the top level is composed of:
\begin{itemize}
    \item Logical Link Control and Adaptation Protocol (L2CAP), takes care of protocol/channel multiplexing with error and flow control
    \item Security Manager, responsible for device pairing and key exchange
    \item Attribute Protocol, allows the device to expose a set of attributes and relative values
    \item Generic Attribute Profile (GATT), defines a way to interact with those attributes: devices can be clients or servers. 
    \item Generic Access Profile (GAP), regulates the procedures of discovery and connection with other devices.
\end{itemize}
\paragraph{Data Exchange.}When the connection completes, a device will assume the role of the GATT client, while the other one would be the GATT server. This means that data exchange happens only at the GATT level: the server will expose a number of \textit{characteristics}, each one being an array of bytes (with the maximum set to 512 B), grouped into \textit{services}. The GATT client can interact with those characteristics, depending on the \textit{properties} of each one. Those properties can be:
\begin{itemize}
    \item \textbf{Read}: the client explicitly requires the value of the characteristic, the server answers accordingly
    \item \textbf{Write}: the client writes the value of the characteristic
    \item \textbf{Notify/Indicate}: the server will inform the client of a new value of the characteristic, without the need of an explicit request. Indications also include an application level acknowledgement of the notification. 
\end{itemize}
Usually notifications are preferred, as only one packet should be exchanged for each new value: no request and no acknowledgment are necessary. 

A particular characteristic of the BLE protocol is the usage of the so-called Connection Events (CE). Data exchange can only happen during these events, and between them the radio goes to sleep: this duty cycled behaviour is one of the reasons of the low energy consumption of the standard. The time interval between two consecutive CEs is called Connection Interval (CI), with allowed values going from 7.5~ms to 4~s (at 1.25~ms steps). These parameters are negotiated at the connection setup, but they can be changed at any time during the connection. 

\paragraph{BLE Throughput.}  The standard defines the data transfer rate to be 1~Mbit/s, but the effective value depends on link quality, layer overheads and hardware constraints. Some studies tried to calculate it, obtaining different results:
\begin{itemize}
    \item Gomez (2011) \cite{gomez2011modeling} reports a maximum throughput of 236.7~kbit/s, by representing the connection through an analytical model. The main flaw of this approach is that it does not take into consideration the hardware constraints of low-cost BLE nodes. 
    \item In \cite{mikhaylov2013performance}, performances of Bluetooth Low Energy are compared to other wireless solutions, reporting a link layer throughput of 122.6~kbit/s: this value is not considering upper layers of the standard, hence it is not considering the packet overhead coming from the application layer.
    \item In \cite{gomez2012overview} presents an overview on the standard, carrying out a practical experiment on throughput: two CC2540 devices, from Texas Instruments, are configured to send data at the maximum possible speed, with Connection Interval set to the minimum possible value. The resulting throughput is 58.48~kbit/s.
    \item \cite{giovanelli2015bluetooth} is probably the most extensive one, as multiple BLE modules are used, each one with its own stack and firmware. It underlines the importance of Connection Intervals in maximizing the throughput. To sum up, the study concluded that the maximum application level throughput that can be achieved is 64~kbit/s, by fine tuning the values of CIs and by using buffer-like structures to maximize the data to send out at each Connection Event.  
\end{itemize}
